\section{Introduction}\label{sec:intro}
The Weil criterion makes the Riemann hypothesis the positivity of an explicit
quadratic form on compactly supported test functions. Localized to an interval
$I_L=(-L/2,L/2)$ it is the form
\begin{equation}\label{eq:target-preview}
 Q_L[f]=K[E_Lf]+w_0\norm f^2+P_L[f]
 -2\!\!\sum_{m\log p<L}\!\!(\log p)p^{-m/2}\re\ip{E_Lf}{U_{m\log p}E_Lf},
\end{equation}
an archimedean energy $K$, a negative contact constant $w_0$, a rank-two form
$P_L$ carrying the two poles of $\zeta$, and a sum over prime powers. The
companion investigation of this program asks for a physical system whose state
norm is \eqref{eq:target-preview}; the present paper asks a prior and entirely
unconditional question, namely what \eqref{eq:target-preview} is.

\subsection{What is proved here}
The four terms of \eqref{eq:target-preview} are not four things. Section
\ref{sec:density} shows that the first two are one: the archimedean symbol
together with the contact is exactly $2\theta'(\tau)$, twice the derivative of
the Riemann--Siegel theta function, so that
\begin{equation}\label{eq:density-preview}
 K[E_Lf]+w_0\norm f^2=\int_\R|\widehat{E_Lf}(\tau)|^2\dd\bar N(\tau),
 \qquad \bar N(T)=\frac{\theta(T)}\pi+1,
\end{equation}
integration against the smooth zero-counting measure. The contact constant then
requires no separate account: it is the $-\log\pi$ inside $\theta'$. Three
consequences follow at once. The archimedean part is a \emph{signed} measure,
negative exactly below $\tau_*=6.2898\ldots$, and once the prime contacts are
added the negative band shrinks with $L$ and is empty above $\log13$; the
crossing at $p=13$ is precisely the point at which the summed prime contacts
overtake $-2\theta'(0)=-w_0$; and the symbol has the graded expansion
$\log(\tau/2\pi)-\tfrac1{24}\tau^{-2}-\tfrac7{960}\tau^{-4}+\cdots$, with no odd
orders, which is four separate conditions on a candidate rather than one.

Section~\ref{sec:jump} shows that the prime sum joins them. Adding and
subtracting the value of the prime symbol at $\tau=0$ puts everything but the
pole term in L\'evy--Khinchine form,
\begin{equation}\label{eq:jump-preview}
 Q_L[f]=\EL[E_Lf]+\gamma_L\norm f^2+P_L[f],
\end{equation}
where $\EL$ is the jump Dirichlet form of the nonnegative L\'evy measure
\begin{equation}\label{eq:measure-preview}
 \mu_L=2\gammak(r)\dd r+2\!\!\sum_{m\log p<L}\!\!(\log p)p^{-m/2}\delta_{m\log p}
\end{equation}
and $\gamma_L$ is an explicit constant. The archimedean density and the prime
atoms are the continuous and atomic parts of one measure: a source for
\eqref{eq:target-preview} is a jump process whose jumps are the primes.

Section~\ref{sec:compression} uses \eqref{eq:density-preview} to sharpen the
compression presentation of the companion manuscript, splitting the form by the
sign of its symbol rather than by a constant, and Section~\ref{sec:cone} takes
up what \eqref{eq:jump-preview} buys in the way of positivity structure. Jump
forms carry Perron--Frobenius structure: their resolvents are positivity
preserving. We verify that here and show that the rank-two pole form destroys
it, at every $L$ and however positive the constant is made --- a third
independent appearance of the pole term as the obstruction of this subject. We
then say precisely what a cone can and cannot do. For a self-adjoint form
Krein--Rutman adds nothing to the spectral theorem; Birkhoff's contraction
bounds a spectral gap and never a spectral radius; and Collatz--Wielandt, which
does bound a spectral radius, needs a pointwise supersolution, which is exactly
what the pole removes.

Section~\ref{sec:necessary} applies the cone to the part of the form that is
still inside one. Writing $Q_L=S_L+P_L$ with $S_L=\EL+\gamma_L\norm\cdot^2$, the
form $S_L$ is of Stieltjes type with constant diagonal, so its least eigenvalue
is a Perron-root deficit $-\alpha_L$ of an explicitly nonnegative kernel. It has
exactly one negative direction, and that direction is $\cosh(x/2)$ to four
digits --- which is the \emph{positive} direction of the pole form. Evaluating
$Q_L$ there gives
\begin{equation}\label{eq:necessary-preview}
 Q_L\geq0\ \hbox{on}\ I_L\ \Longrightarrow\ \alpha_L\leq2\sinh(L/2)+L,
\end{equation}
and because $\alpha_L$ is a Perron root, a single strictly positive trial
function bounds it from below. A trial function violating
\eqref{eq:necessary-preview} would show that the Riemann hypothesis fails. A
scan of $7713$ intervals finds no violation, and finds the least relative margin
to be $9.5\times10^{-4}$.

Section~\ref{sec:selection} collects eight necessary conditions on a source, and
Section~\ref{sec:outlook} says what remains.

\subsection{What is not claimed}
Nothing here establishes Weil positivity on any interval, and no statement
assumes the Riemann hypothesis. The presentations of
Sections~\ref{sec:density}--\ref{sec:compression} are exact rearrangements: they
do not make the problem smaller, and Section~\ref{sec:compression} says so
quantitatively. The necessary condition of Section~\ref{sec:necessary} is
necessary and not sufficient, and Remark~\ref{rem:insufficient} gives the
quantitative reason: it tests the one direction in which the difficulty is
absent. Its holding over the range scanned is the expected outcome and is not
evidence for the Riemann hypothesis. Every numerical statement is labelled as
such, and the check programs of Appendix~\ref{app:checks} state what they do and
do not establish.
